% Version de dépôt (grâce d'édition) — Submission 157 (JURIX 2026, long paper, ≤ 10 pages hors références).
% Format IOS Press officiel. Protocole RÉEL (15 sept. 2026) : templates extraits automatiquement (Claude Opus 5,
% sans relecture juriste), requêtes v0.2 gelées (hash) avant toute extraction du hold-out, une passe d'extraction
% sur le hold-out, stabilité mesurée sur le pilote de conception. Tableaux T3/T4 générés par
% legal-kg/src/reporting/make_tables.py dans tables/. Chaque \todo{} doit disparaître avant dépôt.
\documentclass{IOS-Book-Article}
\usepackage{mathptmx}
\usepackage[T1]{fontenc}
\usepackage[utf8]{inputenc}
\usepackage{graphicx}
\usepackage{booktabs}
\usepackage{array}
\usepackage{amsmath}
\newcommand{\corpus}{CLAUDETTE}
\newcommand{\gl}[1]{item~(\textit{#1})}
\newcommand{\gi}[1]{(\textit{#1})}
\newcommand{\todo}[1]{\textbf{[TODO: #1]}}

\begin{document}
\begin{frontmatter}
\title{Executing the Grey List: Interpretable Queries for Unfair-Term Detection in Terms of Service}
\runningtitle{Executing the Grey List}
\author[A]{\fnms{Ahmed} \snm{El Azhar Jebbari}%
\thanks{Corresponding Author: Ahmed El Azhar Jebbari, LORIA, Campus Scientifique, 54506 Vandœuvre-lès-Nancy, France; e-mail: jebbariazhar@gmail.com.}},
\author[A]{\fnms{Jean-Charles} \snm{Lamirel}},
\author[B]{\fnms{Fatima Zahra} \snm{Boulaich}}
and
\author[C]{\fnms{Fatima} \snm{Ouali}}
\runningauthor{A. El Azhar Jebbari et al.}
\address[A]{LORIA, Université de Lorraine, CNRS, Inria, Nancy, France}
\address[B]{Juristiadata, France}
\address[C]{Faculty of Legal Sciences, Rabat, Morocco}

\begin{abstract}
The Annex of Directive 93/13/EEC lists the types of contractual terms that may be regarded as unfair: an indicative, non-exhaustive list commonly called the grey list. Automated detectors ignore it and classify sentences with learned models whose decisions cite no legal ground. We turn the grey list into an executable resource. Each clause of a contract, identified by its human-validated theme, is mapped to a simple legal template describing who acts, what is obliged, permitted or forbidden, and under which conditions, deadlines and remedies; templates are extracted by a language model that sees neither the queries nor the reference labels. Unfairness queries then translate grey-list items into conditions over these templates; they were frozen, with a recorded hash, before any computation on held-out contracts. Because the list is indicative rather than binding, how well it actually retrieves unfair clauses is an open empirical question, and this is precisely what we measure. On held-out terms of service, we report which grey-list items can be expressed this way, the precision and recall of each query against reference unfairness labels, whether the templates add signal beyond the clause theme alone, and what this interpretability costs compared with a learned text-only model. Every flag comes with the template fields and sentences that triggered it, and false positives are analysed. Queries, templates and code are released.
\end{abstract}

\begin{keyword}
unfair contract terms\sep terms of service\sep Directive 93/13\sep
legal knowledge representation\sep explainable AI\sep information extraction
\end{keyword}
\end{frontmatter}

\section{Introduction}
\label{sec:intro}
Directive 93/13/EEC on unfair terms in consumer contracts~\cite{eu1993directive} defines unfairness by a general test---a significant imbalance in the parties' rights and obligations, contrary to good faith---and complements it with an Annex: an indicative and non-exhaustive list of seventeen kinds of terms, \gl{a} to \gl{q}, that \emph{may} be regarded as unfair. Practitioners call it the grey list. It is the closest thing European consumer law offers to an operational definition of the phenomenon, and it is written in terms of \emph{who} may do \emph{what} to \emph{whom} under \emph{which} conditions: a seller who may terminate a contract of indeterminate duration without reasonable notice \gl{g}, alter the terms unilaterally without a valid reason specified in the contract \gl{j}, or exclude liability for personal injury \gl{a}.

Automated detectors of unfair terms do not use it. Since \corpus{}~\cite{lippi2019claudette}, the reference benchmark for online terms of service (ToS), the task has been framed as sentence classification: a learned model assigns one of eight unfairness categories to each sentence, and its decision cites no legal ground. Rationale-based extensions~\cite{ruggeri2022memory,liepina2020claudettetool} attach explanations to predictions, but the explanations are learned too. Large language models prompted on the same task are neither more accurate nor more accountable~\cite{panarelli2025worth,frasheri2024llm}.

This paper asks a simple question: \textbf{what happens if the grey list is executed literally?} We turn each item into a query over a structured representation of the contract: a \emph{clause template} recording who acts, what is obliged, permitted or forbidden, and under which conditions, deadlines and remedies, attached to clauses identified by a human-validated thematic layer over the 50 \corpus{} contracts. Queries are written before any template is extracted and frozen before any computation on held-out contracts, so that the result is a measurement of the grey list and not of our ability to fit it. Three questions structure the paper, each with a hypothesis recorded in advance (Table~\ref{tab:rq}). \textbf{RQ-A} (Section~\ref{sec:results}): which items can be expressed, and do the frozen queries retrieve the reference labels better than the clause theme alone? \textbf{RQ-B} (Section~\ref{sec:cost}): what does this interpretability cost against learned text-only models? \textbf{RQ-C} (Section~\ref{sec:audit}): what do the false positives reveal about the benchmark?

Because the Annex is indicative rather than binding, nothing guarantees that its items, executed literally, retrieve what the \corpus{} annotators marked as unfair. This is exactly what we measure. Our contributions are: (i) an executable, frozen encoding of the grey list, with an analysis of which items can be expressed as template conditions; (ii) a pre-registered, per-item evaluation against reference labels on held-out contracts, including the gain over the clause theme alone; (iii) a cost accounting of interpretability against text-only models trained on the same split; (iv) an error analysis in which every flag is traceable to the template fields and sentences that triggered it; and (v) the release of the queries with their hashes, the extracted templates and the code.

\section{Background and Related Work}
\label{sec:related}
\textbf{Unfair-term detection.} \corpus{}~\cite{lippi2019claudette,lagioia2019deep} labels sentences of 50 ToS with eight categories---arbitration, unilateral change, content removal, jurisdiction, choice of law, limitation of liability, unilateral termination and contract by using---whose definitions are themselves derived from the Annex and from doctrine~\cite{micklitz2017empire}. The corpus has been extended to several languages~\cite{drawzeski2021corpus,galassi2024multilingual}, included in LexGLUE~\cite{chalkidis2022lexglue}, and used to compare fine-tuned encoders with prompted LLMs~\cite{panarelli2025worth}. Related resources annotate clause validity in German consumer contracts~\cite{braun2024agbde} and potentially abusive clauses in Chilean ToS~\cite{loffler2025chilean}. All of them predict; none of them cites the legal ground of a prediction.

\textbf{Legal knowledge representation.} Rule languages such as LegalRuleML~\cite{athan2015legalruleml} and contract specification languages such as Symboleo~\cite{sharifi2020symboleo} formalise norms as deontic statements with parties, actions and conditions; they are designed for authoring and reasoning, not for being populated from natural-language contracts at scale. Deontic modality extraction~\cite{chalkidis2018obligation,sancheti2022lexdemod} recovers obligations, permissions and prohibitions from text, with macro-F1 in the 0.6--0.8 range depending on the setting---good enough to be an input signal, not good enough to be a ground truth. Our templates sit between these two traditions: a fixed, shallow structure that lawyers can validate clause by clause and that queries can be written against.

\textbf{Queries over contracts.} ContractNLI~\cite{koreeda2021contractnli} evaluates document-level hypotheses (``the receiving party may share confidential information with employees'') by learned entailment; CUAD~\cite{hendrycks2021cuad} retrieves clause types for due diligence. Both pose the question we pose---does this contract contain a provision of this kind?---but answer it with a learned model. We answer it by executing an explicit condition and returning the template fields that satisfy it.

\textbf{Explanation in AI and law.} The field has long argued that legal decision support must explain in legal terms~\cite{atkinson2020explanation}. For unfair-term detection, explanations have so far been rationales selected by attention~\cite{ruggeri2022memory} or generated by a model~\cite{liepina2020claudettetool}. A query that fires because a termination clause grants the provider a discretionary right without notice explains itself by construction, and the explanation can be wrong only in a way a lawyer can check.

\section{The Executable Grey List: from the Annex to Frozen Queries}
\label{sec:greylist}
\label{sec:queries}
Table~\ref{tab:greylist} lists the seventeen items of the Annex, the \corpus{} category each item most directly corresponds to, the clause theme under which such a term is found, and whether the item can be expressed as a condition over the clause templates defined below. Three groups emerge. \emph{Structural items} describe a configuration of actor, modality, action and condition that a template can carry directly: items \gi{a}, \gi{b}, \gi{d}, \gi{f}, \gi{g}, \gi{j}, \gi{k}, \gi{l}, \gi{m}, \gi{p} and \gi{q}. \emph{Quantitative items} require a judgement of proportion or reasonableness that no template field encodes---``disproportionately high compensation'' \gi{e}, ``unreasonably early'' \gi{h}---and can be expressed only with an explicit, declared threshold. \emph{Procedural items} depend on facts outside the text, such as whether the consumer had a real opportunity of becoming acquainted with the terms \gi{i}, and can be approximated at best by a proxy (acceptance by mere use, the \corpus{} category ``contract by using''). Several categories of \corpus{} have no counterpart in the Annex: choice of law is not listed, and jurisdiction clauses are covered only through the general hindrance of legal action in \gi{q}. Conversely, items \gi{c}, \gi{n} and \gi{o} describe terms that are rare in ToS and absent from the reference categories. The grey list and the benchmark therefore overlap without coinciding, and the evaluation must be read item by item.

\begin{table}[t]
\caption{The Annex of Directive 93/13/EEC (grey list), its correspondence with \corpus{} categories (A~arbitration, CH~unilateral change, CR~content removal, J~jurisdiction, LAW~choice of law, LTD~limitation of liability, TER~unilateral termination, USE~contract by using) and with clause themes, and whether the item is expressible as a template condition (S~structural, Q~quantitative with declared threshold, P~procedural proxy, --~not expressible from the text).}
\label{tab:greylist}
\centering
\footnotesize\setlength{\tabcolsep}{3pt}%
\begin{tabular}{@{}c p{5.3cm} p{1.75cm} p{2.7cm} c@{}}
\toprule
Item & Term that has the object or effect of\ldots & \corpus{} & Clause theme & Expr. \\
\midrule
(a) & excluding/limiting liability for death or personal injury & LTD & Limitation of liability & S \\
(b) & excluding/limiting the consumer's legal rights in case of non-performance & LTD & Limitation of liability; Warranty & S \\
(c) & binding the consumer while the seller's performance depends on his will alone & -- & Preamble/scope & -- \\
(d) & retaining sums paid when the consumer withdraws, without reciprocity & -- & Fees and payment & S \\
(e) & requiring disproportionately high compensation from the consumer & -- & Fees and payment & Q \\
(f) & dissolving the contract at the seller's discretion without reciprocity & TER, CR & Termination; User content & S \\
(g) & terminating a contract of indeterminate duration without reasonable notice & TER & Termination & S \\
(h) & automatically extending a fixed-term contract with an unreasonably early deadline & -- & Fees and payment & Q \\
(i) & binding the consumer to terms he had no real opportunity to know & USE & Preamble/scope & P \\
(j) & altering the terms unilaterally without a valid reason specified in the contract & CH & Modification of terms & S \\
(k) & altering unilaterally the characteristics of the product or service & CH & Modification of terms & S \\
(l) & determining the price at delivery or increasing it without a right to cancel & CH & Fees and payment & S \\
(m) & letting the seller decide conformity or interpret any term exclusively & -- & Preamble/scope; Warranty & S \\
(n) & limiting the seller's obligation to respect commitments of his agents & -- & Third-party services & -- \\
(o) & obliging the consumer to perform while the seller does not & -- & -- & -- \\
(p) & transferring the seller's rights and obligations without the consumer's agreement & -- & Preamble/scope & S \\
(q) & excluding or hindering legal action, imposing arbitration, restricting evidence & A, J & Arbitration and disputes & S \\
\bottomrule
\end{tabular}
\end{table}

\textbf{Clause templates.} A template is a flat record attached to a clause (a maximal run of sentences sharing a theme). Its fields are fixed and few: the \emph{actor} (provider, user, third party); the \emph{modality} (obligation, permission, prohibition, power); the \emph{action}, drawn from a closed inventory per theme (terminate, suspend, modify terms, modify service, change price, exclude liability, cap liability, license content, remove content, impose arbitration, waive class action, choose forum, choose law, \ldots); the \emph{object} or affected party; the \emph{conditions} under which the action may be taken (for cause, for a specified reason, at discretion, none stated); the \emph{notice or deadline} (a duration, ``reasonable'', none, or not stated); and the \emph{remedy} left to the other party (refund, right to cancel, compensation, none, not stated). Two conventions matter for the queries. \emph{Not stated} is a value, distinct from \emph{none}: a termination clause that says nothing about notice is not the same as one that excludes it, and grey-list items are sensitive to the difference. And a clause may carry several templates when it states several norms; queries are evaluated per template and aggregated per clause.

\textbf{Queries.} A query is a conjunction of conditions over template fields, optionally over several templates of the same document, that returns the templates satisfying it together with the sentences they were extracted from. The queries were written by the authors from the text of the Annex and the \corpus{} category definitions, before any template was extracted.\todo{disclose AI assistance in drafting the queries if the venue requires it} They were revised once, on development contracts only: six action codes were added to the inventories, and two queries were extended, Q(q) to contractual limitation periods and Q(i) to acceptance deemed by continued use. The set was then frozen with a recorded hash, before any held-out template was extracted, and never edited afterwards. Every query cites the grey-list item it implements. Four examples, in prose:

\begin{itemize}
\item \textbf{Q(g)} --- theme \emph{termination}; actor \emph{provider}; modality \emph{permission} or \emph{power}; action \emph{terminate} or \emph{suspend}; conditions $\in$ \{at discretion, none stated\}; notice $\in$ \{none, not stated\}.
\item \textbf{Q(j)} --- theme \emph{modification of terms}; actor \emph{provider}; action \emph{modify terms}; conditions $\in$ \{at discretion, none stated\} and no valid reason specified; a variant Q(j$'$) additionally requires remedy $\in$ \{none, not stated\}, that is, no right to cancel.
\item \textbf{Q(a/b)} --- theme \emph{limitation of liability}; action \emph{exclude liability} or \emph{cap liability}; object includes personal injury or gross negligence, or the exclusion is unbounded (no carve-out for mandatory law).
\item \textbf{Q(q)} --- action $\in$ \{impose arbitration, waive class action, restrict evidence, shift burden of proof, waive jury trial, limit the claim period\} with modality \emph{obligation} or \emph{prohibition} on the user; or the provider holding the \emph{power} to impose arbitration or choose the forum.
\end{itemize}

Queries come in three families: \emph{presence} of a forbidden configuration (the four above); \emph{absence} of an element the theme's schema expects (a modification clause with no notification mechanism and no notice); and \emph{composition} across clauses of the same document (a discretionary termination right of the provider with no symmetric termination right for the user). All are scored at the sentence level through the sentences their templates cite. A further composition pattern, combining an exclusion of liability, a warranty disclaimer and an indemnification duty of the user, is reported qualitatively only.

\section{Materials: Corpus, Thematic Layer and Clause Templates}
\label{sec:materials}
\label{sec:layer}
\label{sec:templates}
\textbf{Corpus.} \corpus{} comprises 50 online terms of service split into 9,414 sentences, of which 1,032 (11.0\,\%) carry at least one of 1,137 unfairness labels. The 17 held-out contracts contain 4,078 sentences, 440 of them labelled unfair.

\textbf{Thematic layer.} The clause themes we rely on come from a companion resource: a multi-label thematic annotation of the 9,414 sentences of \corpus{} by three trained annotators, with $\alpha$-MASI 0.658 on the 20-theme taxonomy and 0.725 on its 11-theme consolidation, and a leave-one-annotator-out human ceiling of 0.871 accuracy ($\kappa$ 0.859). Agreement measures follow standard practice~\cite{krippendorff2018content,passonneau2006masi,artstein2008intercoder}; disagreement is preserved rather than erased~\cite{braun2024differ}. Every sentence of the 50 contracts carries a primary clause theme and optional secondary themes from a closed vocabulary of 20 themes, assigned independently by three trained annotators and resolved by a cascade (strict agreement 46.8\,\%, majority 48.3\,\%, human arbitration 4.9\,\%). Themes are not unfairness labels, but they are far from independent of them: Table~\ref{tab:themes} gives, for each theme, the share of sentences carrying a \corpus{} unfairness label. Four themes---governing law, modification of terms, termination and limitation of liability---concentrate the unfairness signal with lifts between 3.4 and 5.0 over the 11.0\,\% base rate; the framing themes (preamble, boilerplate, communications) are almost never unfair. This stratification matters twice: it is the reason why the theme alone is a strong baseline that any structured method must beat, and it is the reason why the 11-theme consolidation used downstream never merges themes across strata.

\begin{table}[t]
\caption{Unfairness by clause theme on the 50 \corpus{} contracts (consensus theme; $n$ sentences, share carrying a \corpus{} label, lift over the 11.0\,\% base rate). Themes with fewer than 20 sentences are omitted.}
\label{tab:themes}
\centering
\footnotesize\setlength{\tabcolsep}{3.5pt}%
\begin{tabular}{@{}lrrr@{\quad}lrrr@{}}
\toprule
Theme & $n$ & unfair & lift & Theme & $n$ & unfair & lift \\
\midrule
Governing law & 164 & 54.3\,\% & 4.95 & Fees and payment & 1,077 & 5.5\,\% & 0.50 \\
Modification of terms & 392 & 49.0\,\% & 4.47 & Warranty disclaimer & 613 & 5.4\,\% & 0.49 \\
Termination & 422 & 43.4\,\% & 3.96 & Eligibility and account & 761 & 4.2\,\% & 0.38 \\
Limitation of liability & 604 & 37.2\,\% & 3.40 & Privacy and data & 278 & 4.0\,\% & 0.36 \\
Arbitration and disputes & 847 & 10.5\,\% & 0.96 & Licence and IP & 1,225 & 3.4\,\% & 0.31 \\
User content & 658 & 9.7\,\% & 0.89 & Acceptable use & 911 & 3.4\,\% & 0.31 \\
Preamble and scope & 677 & 9.5\,\% & 0.86 & Miscellaneous boilerplate & 586 & 0.9\,\% & 0.08 \\
Third-party services & 464 & 6.5\,\% & 0.59 & Meta, communications & 184 & 0.0\,\% & 0.00 \\
\bottomrule
\end{tabular}
\end{table}

\textbf{Template extraction.} Templates are extracted automatically; clauses longer than eight sentences are split into windows of at most eight. A language model (Claude Opus~5) receives the clause sentences, the clause theme, the closed action inventory for that theme and the output schema, and returns zero or more templates, each citing the sentences that support it. It runs in an isolated context that contains neither the queries nor any reference label. Outputs are checked for schema conformance, structural validity (cited sentences exist, actions belong to the inventory) and semantic coherence; a lexical anchoring check then resets to \emph{not stated} any condition, notice or remedy value that no cue in the cited sentences supports, a deliberately conservative correction. Templates were not reviewed by legal experts in this study (Section~\ref{sec:discussion}). On a development pilot of 100 clauses extracted three times, 98\,\% of outputs conformed to the schema, the decisive fields (actor, modality, action, condition, notice, remedy) agreed across runs with a mean Jaccard index of 0.81, and the sentences flagged by the frozen queries agreed across runs with Fleiss' $\kappa = 0.86$. The held-out contracts were extracted in a single pass.

\section{Pre-registered Evaluation Protocol}
\label{sec:protocol}
\begin{table}[t]
\caption{Pre-registered questions, hypotheses and tests, recorded before any computation on the held-out contracts. Components requiring expert judgement are deferred to the follow-up study.}
\label{tab:rq}
\centering
\footnotesize\setlength{\tabcolsep}{3pt}%
\begin{tabular}{@{}p{0.55cm} p{5.9cm} p{2.3cm} p{2.3cm}@{}}
\toprule
RQ & Hypothesis H1 & Metric & Test \\
\midrule
A & precision $\geq 0.60$ on the mapped categories; recall $\geq 0.50$ on LTD, TER and CH; $F_1$ gain over theme only on these three & per item and query: P, R, $F_1$ & document-level bootstrap CI \\
B & loss in $F_1$ $\leq 0.05$ against the best text-only model (explanation sufficiency: deferred) & $F_1$ difference & non-inferiority, one-sided CI \\
C & $\geq 20\,\%$ of false positives are grounded flags missing from the reference (expert judgement: deferred) & false positives by item, theme, severity & descriptive only \\
\bottomrule
\end{tabular}
\end{table}

\textbf{Split.} The 50 contracts are split by document into 33 development contracts, used to build the field inventories, to debug the extraction and to revise the queries once, and 17 held-out contracts on which all numbers are reported. \textbf{Pre-registration.} The template schema and the queries were frozen, with a recorded hash and date, before the held-out templates were extracted; the extraction prompt is versioned and hashed as well. One change to the originally planned protocol was recorded at the same time: templates are used as extracted, without expert review, so the pipeline is evaluated end to end. \textbf{Reference.} A sentence is positive for a \corpus{} category if it carries that label in the original corpus; a query is scored against the category or categories its grey-list item maps to (Table~\ref{tab:greylist}), and, for items without a \corpus{} counterpart, against any unfairness label. \textbf{Metrics.} A flagged template marks only the sentences it cites. Per query: precision, recall and $F_1$ at the sentence level; per item: the same, aggregating the item's queries; overall: micro- and macro-averaged over items, with document-level bootstrap intervals; and, for comparison with learned models, a binary \emph{any unfairness} $F_1$. \textbf{Baselines.} (i) \emph{Theme only}: the best achievable retrieval from the clause theme alone, $\mathrm{P}(\text{unfair} \mid \text{theme})$ estimated leave-one-document-out; (ii) a lexical floor, TF-IDF with logistic regression; (iii) a fine-tuned legal encoder~\cite{chalkidis2020legalbert}. Both learned models are trained on the development contracts and scored on the same held-out sentences; the difference in $F_1$ is the price paid for decisions that cite a legal ground. \textbf{Error analysis.} False positives and false negatives are compared with the reference labels and categorised by item, theme and \corpus{} severity, without expert adjudication.

\section{Results A: Expressibility and Per-Item Retrieval}
\label{sec:results}
\textbf{The theme alone is a strong baseline.} On the full corpus, the theme-only predictor reaches an average precision of 0.404 against consensus themes and 0.418 against the arbitrated gold (base rate 0.110), that is, 3.7 times the chance level; the identity of the \emph{combination} of themes carried by a clause reaches 0.589 in document-grouped cross-validation. Any structured method that does not clearly exceed these figures adds explanation but not signal.

\textbf{Expressibility.} Eleven of the seventeen items are structural and expressible without a threshold; two require a declared threshold; one is a procedural proxy; three cannot be expressed from the text of the clause (Table~\ref{tab:greylist}). The items that carry most of the \corpus{} labels---\gi{a}/\gi{b}, \gi{f}/\gi{g}, \gi{j}/\gi{k} and \gi{q}---are all in the structural group.

\textbf{Per-query retrieval.} Table~\ref{tab:queries} reports, for each frozen query on the held-out contracts, the number of flagged sentences, precision and recall against the mapped reference category, and the gain over the theme-only baseline. \todo{one-sentence summary from E2.md: overall precision and recall, best items, gain over theme only}

\IfFileExists{tables/T3_queries.tex}{\begin{table}[t]
\caption{Held-out evaluation of the frozen queries against CLAUDETTE labels.}
\label{tab:queries}
\centering
\footnotesize\setlength{\tabcolsep}{3.5pt}
\begin{tabular}{@{}llrrrrrr@{}}
\toprule
Query & Item (ref.) & $n_+$ & Flagged & P & R & $F_1$ [95\% CI] & $\Delta F_1$ vs theme \\
\midrule
Q-g & (g) TER & 102 & 40 & 0.68 & 0.26 & 0.38 [0.27;0.49] & \\
Q-f-asymmetry & (f) CR,TER & 140 & 12 & 0.83 & 0.07 & 0.13 [0.03;0.26] & \\
Q-j & (j) CH & 88 & 41 & 0.63 & 0.30 & 0.40 [0.32;0.50] & \\
Q-j-strict & (j) CH & 88 & 19 & 0.68 & 0.15 & 0.24 [0.15;0.36] & \\
Q-j-absence & (j) CH & 88 & 8 & 0.62 & 0.06 & 0.10 [0.02;0.23] & \\
Q-k & (k) CH & 88 & 22 & 0.50 & 0.12 & 0.20 [0.08;0.34] & \\
Q-l & (l) CH & 88 & 18 & 0.61 & 0.12 & 0.21 [0.13;0.29] & \\
Q-ab & (a,b) LTD & 132 & 23 & 0.74 & 0.13 & 0.22 [0.10;0.33] & \\
Q-q & (q) A,J & 45 & 103 & 0.17 & 0.38 & 0.23 [0.15;0.32] & \\
Q-i & (i) USE & 44 & 37 & 0.43 & 0.36 & 0.40 [0.31;0.46] & \\
Q-m & (m) & --- & 39 & \multicolumn{4}{l}{audit only (no reference category)} \\
Q-p & (p) & --- & 17 & \multicolumn{4}{l}{audit only (no reference category)} \\
Q-d & (d) & --- & 21 & \multicolumn{4}{l}{audit only (no reference category)} \\
Q-e & (e) & --- & 0 & \multicolumn{4}{l}{audit only (no reference category)} \\
Q-h & (h) & --- & 0 & \multicolumn{4}{l}{audit only (no reference category)} \\
\midrule
\emph{item (a)} & LTD & 132 & 4 & 0.50 & 0.02 & 0.03 [0.00;0.08] & -0.41 [-0.50;-0.29] \\
\emph{item (b)} & LTD & 132 & 19 & 0.79 & 0.11 & 0.20 [0.08;0.31] & -0.24 [-0.40;-0.08] \\
\emph{item (f)} & TER,CR & 140 & 12 & 0.83 & 0.07 & 0.13 [0.03;0.26] & -0.21 [-0.33;-0.08] \\
\emph{item (g)} & TER & 102 & 40 & 0.68 & 0.26 & 0.38 [0.27;0.49] & +0.02 [-0.10;0.13] \\
\emph{item (i)} & USE & 44 & 37 & 0.43 & 0.36 & 0.40 [0.31;0.46] & +0.21 [0.10;0.30] \\
\emph{item (j)} & CH & 88 & 42 & 0.62 & 0.30 & 0.40 [0.32;0.50] & +0.00 [-0.10;0.12] \\
\emph{item (k)} & CH & 88 & 22 & 0.50 & 0.12 & 0.20 [0.08;0.34] & -0.20 [-0.28;-0.11] \\
\emph{item (l)} & CH & 88 & 18 & 0.61 & 0.12 & 0.21 [0.13;0.29] & -0.19 [-0.33;-0.05] \\
\emph{item (q)} & A,J & 45 & 103 & 0.17 & 0.38 & 0.23 [0.15;0.32] & +0.05 [-0.04;0.14] \\
\midrule
\textbf{Structural items (8)} & & & & 0.46 & 0.15 & 0.22 (macro 0.22) & -0.15 \\
\bottomrule
\end{tabular}
\end{table}
}{%
\begin{table}[t]
\caption{Held-out evaluation of the frozen queries against \corpus{} labels (pending: generated table not found).}
\label{tab:queries}
\centering
\footnotesize\setlength{\tabcolsep}{4pt}%
\begin{tabular}{llccccc}
\toprule
Query & Item & Reference & Flagged & Precision & Recall & $\Delta F_1$ vs theme only \\
\midrule
Q(g)   & (g)   & TER     & --- & --- & --- & --- \\
Q(f)   & (f)   & TER, CR & --- & --- & --- & --- \\
Q(j)   & (j)   & CH      & --- & --- & --- & --- \\
Q(j$'$)& (j)   & CH      & --- & --- & --- & --- \\
Q(k)   & (k)   & CH      & --- & --- & --- & --- \\
Q(l)   & (l)   & CH      & --- & --- & --- & --- \\
Q(a/b) & (a),(b) & LTD   & --- & --- & --- & --- \\
Q(q)   & (q)   & A, J    & --- & --- & --- & --- \\
Q(i)   & (i)   & USE     & --- & --- & --- & --- \\
Q(m)   & (m)   & any     & --- & --- & --- & --- \\
Q(p)   & (p)   & any     & --- & --- & --- & --- \\
\midrule
All structural items & & & --- & --- & --- & --- \\
\bottomrule
\end{tabular}
\end{table}}

\section{Results B: The Price of Interpretability}
\label{sec:cost}
On the same held-out sentences, the fine-tuned legal encoder reaches a binary $F_1$ of 0.701 [0.650; 0.749] and the lexical floor 0.529 [0.492; 0.572]. \todo{binary F1 of the union of queries and the gap, from Table~T4} Table~\ref{tab:cost} gives the full comparison. \todo{non-inferiority test at 0.05 from STATS.md: H1 retained or rejected}

\IfFileExists{tables/T4_cost.tex}{\begin{table}[t]
\caption{The price of interpretability on the held-out sentences (any \corpus{} label, 440 positives among 4,078). All systems are scored by the same evaluation run; the learned models are trained on the 33 development contracts.}
\label{tab:cost}
\centering
\footnotesize\setlength{\tabcolsep}{4pt}%
\begin{tabular}{@{}lrrlrl@{}}
\toprule
System & P & R & $F_1$ [95\,\% CI] & Flagged & Cites a legal ground \\
\midrule
Theme only (leave-one-document-out) & 0.42 & 0.46 & 0.44 & 489 & no \\
TF-IDF + logistic regression & 0.45 & 0.63 & 0.53 [0.49; 0.57] & 614 & no \\
Legal-BERT, fine-tuned & 0.64 & 0.77 & 0.70 [0.65; 0.75] & 527 & no \\
\textbf{Frozen queries (union)} & 0.41 & 0.33 & 0.36 [0.32; 0.42] & 356 & \textbf{yes} \\
\bottomrule
\end{tabular}
\end{table}
}{\begin{table}[t]\caption{The price of interpretability (pending: generated table not found).}\label{tab:cost}\centering\footnotesize ---\end{table}}

\section{Results C: When a Query Fires and the Benchmark Does Not}
\label{sec:audit}
Every flag returns the template fields that satisfied the query and the sentences they were extracted from, so each error can be traced to one of three sources: a mis-filled field, a query broader than its grey-list item, or a missing reference label for a clause that a literal reading of the Annex captures. In this study the attribution is automatic and partial. \todo{from DISAGREEMENT.md: false positives by item, theme and severity; false-negative typology}. Deciding whether a false positive reveals an omission in the reference corpus requires legal judgement; it is left to the expert study described in Section~\ref{sec:discussion}.

\section{Discussion and Limitations}
\label{sec:discussion}
Three points follow. First, the grey list is executable for the items that matter most in ToS, and the execution is cheap once themes and templates exist: a query is a few conditions that a lawyer reads in one line. Second, the theme carries much of the retrievable signal; what the templates add is not primarily recall but \emph{grounds}---the ability to say which item fired and why---and, on some items, precision.\todo{check this claim against E2: gain over theme only per item} Third, the mismatch between the Annex and the benchmark is itself a finding: categories such as choice of law have no grey-list ground, and grey-list items such as \gi{d} or \gi{h} have no reference labels, so a detector that executed the law faithfully would be penalised by the benchmark for being right. \textbf{Limitations and future work.} The templates were not reviewed by legal experts, as this was not feasible within the time frame of this study. Our results therefore characterise the end-to-end pipeline, with extraction treated as a measured source of noise, and they do not establish the correctness of individual templates. The choice of extractor matters: on the development pilot, with the previous prompt version, a second model (gpt-6-astra) agreed with the one used here on the decisive fields with a mean Jaccard index of about 0.6. The held-out set was extracted once. The reference labels are sentence-level while the Annex is clause-level; thresholds for quantitative items are declared, not derived; the held-out set is small and the intervals are wide; the queries reflect one reading of one legal system's list. This work is part of an ongoing collaboration with legal experts, who \todo{state their actual contribution to this study}. The follow-up study addresses these limitations: independent double validation of the templates by two legal experts, with per-field agreement, which will separate extraction errors from query errors; re-extraction with pinned model versions and replication with an open-weight model; and an expert audit of the flagged sentences and of the explanations the queries return.

\section{Conclusion and Availability}
\label{sec:conclusion}
We have treated the grey list of Directive 93/13/EEC as what it is---a list of configurations of rights and obligations---and executed it as frozen queries over clause templates extracted automatically from the \corpus{} corpus. The result is a detector whose every decision cites the item it implements, an item-by-item measurement of how much of the annotated unfairness a literal reading of the Annex retrieves, and an error analysis that turns false positives into questions for the reference corpus. The queries with their hashes, the extracted templates, the code and the annotation platform used to build the thematic layer are released.

\bibliographystyle{vancouver}
\bibliography{references}
\end{document}
